\begin{scanpage}{0097}{80}
\[\begin{aligned}\left.\begin{matrix}\overline\Sigma^0(\vartheta)\\\Sigma^0(\vartheta)\end{matrix}\right\}={}&\mp\frac{\sin\vartheta}{2\cos\vartheta}(e^{ikl\cos\vartheta}-e^{-ikl\cos\vartheta})[J_1(\zeta)-iJ_2(\zeta)]\\&\mp\frac{\cos\vartheta}{\sin\vartheta}J_1(\zeta)e^{ikl\cos\vartheta}\quad(\zeta=2ka\sin\vartheta).
\end{aligned}\tag{13.29}\]
These equations allow one to determine in the physical optics approach the effective scattering area of a finite cylinder.
\sectiontitle{14}{The Field Created by the Nonuniform Part of the Current}
Let us find the field from the nonuniform part of the current caused by the surface's discontinuity. Figuratively speaking, the field scattered by the cylinder is created by the ``luminous'' regions on its end and lateral surface. Mathematically this field is described by the sum of spherical waves from the ``luminous'' points 1, 2 and 3 (see Figure 24). Obviously the field from the nonuniform part of the current also will have the form of spherical waves diverging from these same points.

In the case when the length and diameter of the cylinder are sufficiently large in comparison with the wavelength, one may approximately consider that the nonuniform part of the current near the discontinuity is the same as that on a corresponding wedge. The field radiated by this part of the current in principle may be found in the same way as in the case of the disk. However, such a method is rather complicated. We will find the desired field by a simpler and more graphic method, starting from a physical analysis of the solution obtained for the disk.

For this purpose, let us investigate the structure of waves (12.01) and (12.02) which are radiated by the disk. These equations include the factor
\[\frac{ia}{\sqrt{2\pi ka(\sin\gamma+\sin\vartheta)}}\frac{e^{ikR}}R=i\frac{e^{ikR}}{\sqrt{2\pi kR}}\sqrt{\frac aR\frac1{\sin\gamma+\sin\vartheta}}.\tag{14.01}\]
\end{scanpage}
